\chapter{Complexity — Uniform Tableau Emitter Instr}
%%% generated by the blueprint pipeline from TCSlib.Complexity.CircuitComplexity.UniformTableauEmitterInstr
\section{Overview}

Declarations of the Lean module \texttt{TCSlib.Complexity.CircuitComplexity.UniformTableauEmitterInstr}.

\section{Declarations}

\begin{theorem}[Updating the circuit-input count]
\lean{Complexity.UTab.ER.set_nc}
\label{Complexity.UTab.ER.set_nc}
\leanok
\uses{Complexity.UTab.ER, Complexity.UTab.ER.set, Complexity.UTab.rNC}
\difficulty{1}
For any register state $e$ and natural number $v$, setting its circuit-input-count
coordinate to $v$ yields the state obtained by replacing that coordinate of $e$ with
$v$.
\end{theorem}

\begin{theorem}[Setting the hard-wired-bit count]
\lean{Complexity.UTab.ER.set_z}
\label{Complexity.UTab.ER.set_z}
\leanok
\uses{Complexity.UTab.ER, Complexity.UTab.ER.set, Complexity.UTab.rZ}
\difficulty{1}
Let $e$ be a record with natural-number fields for the circuit-input count,
hard-wired-bit count, step budget, current instruction, preceding cell, first
instruction of the current step, two loop counters, hard-wired and free input positions,
and scratch space. Setting its hard-wired-bit count to $v$ gives the record obtained by
replacing that field with $v$ and leaving every other field unchanged.
\end{theorem}

\begin{theorem}[Setting the step budget]
\lean{Complexity.UTab.ER.set_t}
\label{Complexity.UTab.ER.set_t}
\leanok
\uses{Complexity.UTab.ER, Complexity.UTab.ER.set, Complexity.UTab.rT}
\difficulty{1}
Let $e$ be a record with natural-number fields for the circuit-input count,
hard-wired-bit count, step budget, current instruction, preceding cell, first
instruction of the current step, two loop counters, hard-wired and free input positions,
and scratch space. Setting its step budget to $v$ gives the record obtained by replacing
that field with $v$ and leaving every other field unchanged.
\end{theorem}

\begin{theorem}[Updating the Current Instruction]
\lean{Complexity.UTab.ER.set_i}
\label{Complexity.UTab.ER.set_i}
\leanok
\uses{Complexity.UTab.ER, Complexity.UTab.ER.set, Complexity.UTab.rI}
\difficulty{1}
For every register state $e$ and natural number $v$, setting coordinate $3$ of $e$ to
$v$ gives the state obtained by replacing its current instruction with $v$.
\end{theorem}

\begin{theorem}[Setting the preceding cell]
\lean{Complexity.UTab.ER.set_q}
\label{Complexity.UTab.ER.set_q}
\leanok
\uses{Complexity.UTab.ER, Complexity.UTab.ER.set, Complexity.UTab.rQ}
\difficulty{1}
Let $e$ be a record with natural-number fields for the circuit-input count,
hard-wired-bit count, step budget, current instruction, preceding cell, first
instruction of the current step, two loop counters, hard-wired and free input positions,
and scratch space. Setting its preceding-cell field to $v$ gives the record obtained by
replacing that field with $v$ and leaving every other field unchanged.
\end{theorem}

\begin{theorem}[Updating the Step-Loop Counter]
\lean{Complexity.UTab.ER.set_lb}
\label{Complexity.UTab.ER.set_lb}
\leanok
\uses{Complexity.UTab.ER, Complexity.UTab.ER.set, Complexity.UTab.rLB}
\difficulty{1}
For any register state $e$ and natural number $v$, updating its step-loop counter to $v$
gives the state that agrees with $e$ in every other register and has step-loop counter
$v$.
\end{theorem}

\begin{theorem}[Updating the inner loop counter]
\lean{Complexity.UTab.ER.set_lt}
\label{Complexity.UTab.ER.set_lt}
\leanok
\uses{Complexity.UTab.ER, Complexity.UTab.ER.set, Complexity.UTab.rLT}
\difficulty{1}
For every register state $e$ and natural number $v$, setting coordinate $6$ of $e$ to
$v$ gives the state obtained by replacing its inner loop counter with $v$ and leaving
all other fields unchanged.
\end{theorem}

\begin{theorem}[Updating the Inner-Loop Counter]
\lean{Complexity.UTab.ER.set_ll}
\label{Complexity.UTab.ER.set_ll}
\leanok
\uses{Complexity.UTab.ER, Complexity.UTab.ER.set, Complexity.UTab.rLL}
\difficulty{1}
For any register state $e$ and natural number $v$, setting the inner-loop counter to $v$
yields the state obtained by replacing its inner-loop counter with $v$.
\end{theorem}

\begin{theorem}[Setting the hard-wired input position]
\lean{Complexity.UTab.ER.set_pa}
\label{Complexity.UTab.ER.set_pa}
\leanok
\uses{Complexity.UTab.ER, Complexity.UTab.ER.set, Complexity.UTab.rPA}
\difficulty{1}
Let $e$ be a record with natural-number fields for the circuit-input count,
hard-wired-bit count, step budget, current instruction, preceding cell, first
instruction of the current step, two loop counters, hard-wired and free input positions,
and scratch space. Setting its hard-wired input position to $v$ gives the record
obtained by replacing that field with $v$ and leaving every other field unchanged.
\end{theorem}

\begin{theorem}[Setting the free input position]
\lean{Complexity.UTab.ER.set_pb}
\label{Complexity.UTab.ER.set_pb}
\leanok
\uses{Complexity.UTab.ER, Complexity.UTab.ER.set, Complexity.UTab.rPB}
\difficulty{1}
Let $e$ be a record with natural-number fields for the circuit-input count,
hard-wired-bit count, step budget, current instruction, preceding cell, first
instruction of the current step, two loop counters, hard-wired and free input positions,
and scratch space. Setting its free input position to $v$ gives the record obtained by
replacing that field with $v$ and leaving every other field unchanged.
\end{theorem}

\begin{theorem}[Setting the scratch field]
\lean{Complexity.UTab.ER.set_tmp}
\label{Complexity.UTab.ER.set_tmp}
\leanok
\uses{Complexity.UTab.ER, Complexity.UTab.ER.set, Complexity.UTab.rTMP}
\difficulty{1}
Let $e$ be a record with natural-number fields for the circuit-input count,
hard-wired-bit count, step budget, current instruction, preceding cell, first
instruction of the current step, two loop counters, hard-wired and free input positions,
and scratch space. Setting its scratch field to $v$ gives the record obtained by
replacing that field with $v$ and leaving every other field unchanged.
\end{theorem}

\begin{theorem}[Execution of a template program]
\lean{Complexity.UTab.goes_tp}
\label{Complexity.UTab.goes_tp}
\leanok
\uses{Complexity.CounterProg.Goes, Complexity.CounterProg.MOp.exec, Complexity.CounterProg.goes_tmpl, Complexity.UTab.Lb, Complexity.UTab.Lmax, Complexity.UTab.Tpl, Complexity.UTab.Tpl.next, Complexity.UTab.length_tmpl_le, Complexity.UTab.prog, Complexity.UTab.tmpl, Complexity.UTab.tpL}
\difficulty{4}
Let $M$ be a finite Turing machine, let $x$ be a binary input, let $\theta$ be a
template, let $\rho$ be a valuation of the eleven registers, and let $p\in\mathbb N$ be
the input position. Starting at the program label for $\theta$, the counter program
reaches the label following $\theta$, with registers $\rho$ and input position $p$
unchanged, and appends the output produced by executing the template's operations with
register values $\rho$. This execution takes at most the number of operations in the
template plus one steps.
\end{theorem}

\begin{theorem}[Execution of a template instruction]
\lean{Complexity.UTab.goes_instr}
\label{Complexity.UTab.goes_instr}
\leanok
\uses{Complexity.CounterProg.Goes, Complexity.CounterProg.Goes.mono, Complexity.CounterProg.Goes.trans, Complexity.CounterProg.MOp.exec, Complexity.CounterProg.goes_goto, Complexity.UTab.ER, Complexity.UTab.ER.f, Complexity.UTab.ER.f_i, Complexity.UTab.ER.f_q, Complexity.UTab.ER.set_i, Complexity.UTab.ER.set_q, Complexity.UTab.Lb, Complexity.UTab.Lmax, Complexity.UTab.Tpl, Complexity.UTab.Tpl.after, Complexity.UTab.Tpl.next, Complexity.UTab.Tpl.t0, Complexity.UTab.goes_incR, Complexity.UTab.goes_tp, Complexity.UTab.prog, Complexity.UTab.rI, Complexity.UTab.rQ, Complexity.UTab.tmpl, Complexity.UTab.tpL}
\difficulty{4}
Let $M$ be a finite-state Turing machine with Boolean input and $k$ work tapes, let $x$
be a Boolean word, let $\theta$ be a template for $M$, let $e$ be a record of eleven
natural-number register values, and let $p\in\mathbb{N}$. Assume that the label assigned
as the successor of $\theta$ is the label that begins its per-template increment phase.
From the label for $\theta$, with registers initialized to those specified by $e$, input
position $p$, and any output word $o$, the counter program reaches the label assigned
after $\theta$ in at most the length of $\theta$'s operation sequence plus three steps,
still at input position $p$. Its registers are those specified by $e$ with the
instruction counter increased by one and, when the template's initial-step flag is
false, the preceding-cell counter also increased by one; its output is $o$ followed by
the bits obtained by executing the template's operations using the initial register
values.
\end{theorem}

\begin{theorem}[Distinct source and destination registers]
\lean{Complexity.UTab.CpS.regs_ne}
\label{Complexity.UTab.CpS.regs_ne}
\leanok
\uses{Complexity.UTab.CpS, Complexity.UTab.CpS.dst, Complexity.UTab.CpS.src, Complexity.UTab.rTMP}
\difficulty{4}
Let $M$ be a finite Turing machine with Boolean tape symbols, let $x$ be a finite
Boolean string, and let $s$ be a copy-phase state for the tape length parameter $M.k$.
Assign source and destination registers to $s$ by the specified copy-phase rules. Then
its source and destination registers are distinct, and neither is register $10$.
\end{theorem}

\begin{theorem}[Register-copying loop]
\lean{Complexity.UTab.goes_copy}
\label{Complexity.UTab.goes_copy}
\leanok
\uses{Complexity.CounterProg.Goes, Complexity.CounterProg.Goes.congr, Complexity.CounterProg.Goes.mono, Complexity.CounterProg.Goes.trans, Complexity.CounterProg.goes_inc, Complexity.CounterProg.goes_loop, Complexity.UTab.CpS, Complexity.UTab.CpS.dst, Complexity.UTab.CpS.exit, Complexity.UTab.CpS.regs_ne, Complexity.UTab.CpS.src, Complexity.UTab.prog, Complexity.UTab.rTMP}
\difficulty{6}
Let $M$ be a finite-state multitape Turing machine with Boolean tape symbols, let $x$ be
a finite Boolean input, and let $s$ be a copy-loop mode. Given eleven-register values
$\rho$ with scratch register value zero, and an input position $p$, the program reaches
the exit mode with the destination register increased by the initial source-register
value, all other register values unchanged, and input position still $p$. For every
output prefix, this occurs within $7a+2$ steps, where $a$ is the initial value of the
source register, and the program leaves the output unchanged.
\end{theorem}

\begin{theorem}[Clearing a Counter Register]
\lean{Complexity.UTab.goes_clear}
\label{Complexity.UTab.goes_clear}
\leanok
\uses{Complexity.CounterProg.Goes, Complexity.CounterProg.Goes.congr, Complexity.CounterProg.goes_loop, Complexity.UTab.ClS, Complexity.UTab.ClS.exit, Complexity.UTab.ClS.reg, Complexity.UTab.Lb, Complexity.UTab.Lmax, Complexity.UTab.goes_refl, Complexity.UTab.prog}
\difficulty{4}
Let $M$ be a multitape Turing machine with Boolean tape symbols, let $x$ be a finite
Boolean input, let $s$ be one of the three clearing modes, let $\rho$ assign a
natural-number value to each of eleven registers, and let $p$ be a natural-number input
position. For every output word $o$, the counter program started at phase $0$ of the
clearing routine for $s$, with registers $\rho$, input position $p$, and output $o$,
reaches the exit for $s$ in at most $2\rho(r)+1$ steps, where $r$ is the register
selected by $s$; at that point only $r$ has changed, its value is $0$, the input
position is still $p$, and the output is still $o$.
\end{theorem}

\begin{theorem}[Encoding a mapped gate list]
\lean{Complexity.UTab.gbits_map}
\label{Complexity.UTab.gbits_map}
\leanok
\uses{BoolCircuit.DAGGate, BoolCircuit.DAGGate.encode, Complexity.UTab.gbits}
\difficulty{1}
Let $l$ be a finite list in a type $\alpha$, and let $f$ map each element of $\alpha$ to
a directed acyclic Boolean gate. The Boolean string formed from the list of gates
obtained by applying $f$ to $l$ is the concatenation, in list order, of a one followed
by the encoding of $f(a)$ for each $a\in l$.
\end{theorem}

\begin{theorem}[Execution of instruction operations]
\lean{Complexity.UTab.instrOps_exec}
\label{Complexity.UTab.instrOps_exec}
\leanok
\uses{BoolCircuit.copyGate, BoolCircuit.shiftGate, BoolCircuit.tabFixed, Complexity.CounterProg.LinE.val, Complexity.CounterProg.MOp.exec, Complexity.UTab.SSrc.toLinE, Complexity.UTab.Tpl, Complexity.UTab.Tpl.kind, Complexity.UTab.Tpl.srcs, Complexity.UTab.baseE, Complexity.UTab.flatMap_exec_sgOps, Complexity.UTab.gbits, Complexity.UTab.gbits_map, Complexity.UTab.instrOps, Complexity.cfgArity, Complexity.cfgF, Complexity.cfgWidth, Complexity.snapWidth}
\difficulty{4}
Let $M$ be a finite-state Turing machine with $k$ work tapes, let $\theta$ be a
template, and let $\rho$ assign a natural number to each of the eleven registers.
Executing the operation list associated with $\theta$ yields the bit string obtained by
encoding, in order, the copy gates for the template’s source expressions evaluated at
$\rho$, followed by the encoding of the fixed gates for its kind shifted by the value of
the base expression at $\rho$.
\end{theorem}

\begin{theorem}[Value of the Base Expression]
\lean{Complexity.UTab.baseE_val}
\label{Complexity.UTab.baseE_val}
\leanok
\uses{BoolCircuit.tabBase, Complexity.CounterProg.LinE.val, Complexity.UTab.ER, Complexity.UTab.ER.f, Complexity.UTab.ER.f_i, Complexity.UTab.ER.f_nc, Complexity.UTab.ER.f_z, Complexity.UTab.baseE, Complexity.cfgArity, Complexity.cfgF, Complexity.cfgWidth}
\difficulty{3}
Let $M$ be a finite Turing machine over the binary alphabet, let $n$ be a natural
number, let $z$ be a finite binary string, and let $e$ be a register state whose
input-length register equals $n$ and whose fixed-string-length register equals $|z|$.
The value at $e$ of the base linear expression for $M$ is $n+2+|z|+W(i+1)$, where $i$ is
the instruction register and $W$ is the table width determined by $M$.
\end{theorem}

\begin{theorem}[Template bits match tableau instruction bits]
\lean{Complexity.UTab.instr_bits}
\label{Complexity.UTab.instr_bits}
\leanok
\uses{BoolCircuit.BitSrc.Valid, BoolCircuit.GInstr, BoolCircuit.padSrc, BoolCircuit.tabInstrGates, Complexity.CfgKind, Complexity.CounterProg.MOp.exec, Complexity.UTab.ER, Complexity.UTab.ER.f, Complexity.UTab.SSrc.Ok, Complexity.UTab.SSrc.toBitSrc, Complexity.UTab.Tpl, Complexity.UTab.Tpl.kind, Complexity.UTab.Tpl.srcs, Complexity.UTab.baseE_val, Complexity.UTab.gbits, Complexity.UTab.gbits_append, Complexity.UTab.instrOps, Complexity.UTab.instrOps_exec, Complexity.UTab.tabSrc_toBitSrc, Complexity.UTab.tmpl, Complexity.cfgArity, Complexity.cfgF, Complexity.cfgWidth, Complexity.snapWidth}
\difficulty{4}
Let $M$ be a finite-state Turing machine with $k$ work tapes, let $n$ be a natural
number, let $z$ be a finite Boolean string, let $\theta$ be a template, and let $e$ be a
register record. Suppose the template operation list equals its instruction-operation
list, the input-count and hard-wired-input-length registers of $e$ are $n$ and $|z|$,
and an instruction has kind equal to the kind of $\theta$. Assume the template’s source
expressions, evaluated at the register record, give exactly the instruction’s source
list, whose length is the configuration arity; every such source is valid for the input
length $|z|+n$, output width determined by $M$, and instruction index in $e$. Then
executing the template produces exactly the encoded bits of the tableau instruction.
\end{theorem}

\begin{theorem}[Validity under source padding]
\lean{Complexity.UTab.ok_padS}
\label{Complexity.UTab.ok_padS}
\leanok
\uses{Complexity.UTab.SSrc, Complexity.UTab.SSrc.Ok, Complexity.UTab.padS}
\difficulty{4}
Let $M$ be a finite-tape Boolean Turing machine, let $m$ be a natural number, let $l$ be
a finite list of sources, and let $\rho$ assign natural-number values to eleven
registers. If every source in $l$ is valid, then every source in the list formed by
appending enough constant false sources to reach length $m$ (or appending none if $l$
already has length at least $m$) is valid.
\end{theorem}

\begin{theorem}[Validity of a conditional previous source]
\lean{Complexity.UTab.ok_nbS}
\label{Complexity.UTab.ok_nbS}
\leanok
\uses{Complexity.UTab.SSrc.Ok, Complexity.UTab.nbS, Complexity.UTab.rQ}
\difficulty{4}
Let $M$ be a finite-state Turing machine over the binary alphabet, let $\rho$ assign a
natural number to each of the eleven registers, let $ok$ be a Boolean, let
$d\in\{0,1,2\}$, and let $j\in\mathbb N$. If $ok$ is true, suppose $1\leq\rho(r_Q)+d$.
Then the source that is the previous-configuration source indexed by $(d,j)$ when $ok$
is true and the constant-false source otherwise is valid for $\rho$.
\end{theorem}

\begin{theorem}[Validity of a conditional source]
\lean{Complexity.UTab.ok_chS}
\label{Complexity.UTab.ok_chS}
\leanok
\uses{Complexity.UTab.SSrc.Ok, Complexity.UTab.chS, Complexity.UTab.rI}
\difficulty{4}
Let $M$ be a finite-tape Boolean Turing machine, let $j$ be a natural number, let $ok$
be a Boolean, and let $\rho$ assign natural numbers to the eleven registers. If $ok$ is
true, assume that $1 \leq \rho(rI)$. Then the source that is the $j$th current-cell
source when $ok$ is true and the constant false source otherwise satisfies the source
validity condition for $\rho$.
\end{theorem}

\begin{theorem}[Validity of Postponed Sources]
\lean{Complexity.UTab.ok_psnS}
\label{Complexity.UTab.ok_psnS}
\leanok
\uses{Complexity.UTab.SSrc.Ok, Complexity.UTab.psnS, Complexity.UTab.rLB}
\difficulty{4}
Let $M$ be a finite-state Turing machine over the binary alphabet, let $\rho$ assign a
natural number to each of the eleven registers, let $W\in\mathbb N$, and let $t_0$ be
Boolean. If $t_0$ is false, assume $1\leq\rho(r_{LB})$. Every source in the length-$W$
list whose entries are constant false when $t_0$ is true and postponed sources at
positions $j=0,\ldots,W-1$ otherwise is valid for $\rho$.
\end{theorem}

\begin{theorem}[Validity of cell sources]
\lean{Complexity.UTab.ok_cellS}
\label{Complexity.UTab.ok_cellS}
\leanok
\uses{Complexity.UTab.SSrc.Ok, Complexity.UTab.cellS, Complexity.UTab.ok_chS, Complexity.UTab.ok_nbS, Complexity.UTab.ok_psnS, Complexity.UTab.rI, Complexity.UTab.rLB, Complexity.UTab.rQ}
\difficulty{4}
Let $M$ be a finite-state Turing machine over the Boolean alphabet, let $\rho$ assign a
natural number to each of eleven registers, let $W\in\mathbb{N}$, and let
$t_0,r_{\mathrm{lo}},r_T,r_{\mathrm{hi}}$ be Boolean values. Suppose that
$\rho(r_Q)\ge1$ whenever $t_0$ is false and $r_{\mathrm{lo}}$ is true, that
$\rho(r_I)\ge1$ whenever $r_{\mathrm{lo}}$ is true, and that $\rho(r_{LB})\ge1$ whenever
$t_0$ is false. Then every source in the cell source list of width $W$ determined by
these Boolean values is valid for $\rho$.
\end{theorem}

\begin{theorem}[Validity of input sources]
\lean{Complexity.UTab.ok_inpS}
\label{Complexity.UTab.ok_inpS}
\leanok
\uses{Complexity.UTab.SSrc.Ok, Complexity.UTab.inpS, Complexity.UTab.layS, Complexity.UTab.ok_chS, Complexity.UTab.ok_nbS, Complexity.UTab.ok_psnS, Complexity.UTab.rI, Complexity.UTab.rLB, Complexity.UTab.rQ}
\difficulty{4}
Let $W$ be a natural number, let $t_0,p_0,p_1,p_{-1}$ be Boolean values, let $
ho$
assign a natural number to each of eleven registers, and let
$
ho_Q,
ho_I,
ho_{LB},
ho_{PA},
ho_Z$ denote its values at registers
$r_Q,r_I,r_{LB},r_{PA},r_Z$. A source is valid when it is a constant, when it is a
zero-input source and $
ho_{PA}<
ho_Z$, when it is a previous-cell source with offset
$d\in\{0,1,2\}$ and $1\leq\rho_Q+d$, when it is a current-cell source and $1\leq\rho_I$,
or when it is a past-position source and $1\leq\rho_{LB}$. Suppose $1\leq\rho_Q$
whenever $t_0$ and $p_0$ are both false, $1\leq\rho_I$ whenever $p_0$ is false, and
$1\leq\rho_{LB}$ whenever $t_0$ is false. If the layer source for a chosen layer in
$\{0,1,2\}$ is valid, then every source in the input list—whose entries include the
specified constant and layer sources, conditional previous-cell and current-cell
sources, and $W$ past-position sources—is valid.
\end{theorem}

\begin{theorem}[Validity of snapshot sources]
\lean{Complexity.UTab.ok_snapS}
\label{Complexity.UTab.ok_snapS}
\leanok
\uses{Complexity.UTab.SSrc.Ok, Complexity.UTab.ok_psnS, Complexity.UTab.rI, Complexity.UTab.rLB, Complexity.UTab.snapS}
\difficulty{4}
Let $M$ be a finite Turing machine with $k$ work tapes, let $W$ be a natural number, let
$t_0$ be a Boolean, and let $\rho$ assign a natural number to each of the eleven
registers. If register $I$ has value at least $1$, and register $LB$ has value at least
$1$ whenever $t_0$ is false, then every source in the snapshot source list of width $W$
and tape count $k$ is valid for $\rho$: its input position is strictly before register
$Z$ when it reads register $PA$, and each referenced register block has a positive
index.
\end{theorem}

\begin{definition}[Instruction gate encoding]
\lean{Complexity.UTab.ibits}
\label{Complexity.UTab.ibits}
\leanok
\uses{BoolCircuit.tabInstrGates, Complexity.UTab.gbits, Complexity.cfgArity, Complexity.cfgF, Complexity.cfgInstr, Complexity.cfgWidth, Complexity.tabLayout}
Let $M$ be a finite binary Turing machine, let $x$ be a finite bit string, and let
$n,z,T,i$ be natural numbers or bit strings as appropriate: the instruction bit string
is the gate encoding of the gates produced by tabulating the instruction selected by the
configuration at index $i$, for input data $z$, parameters $n,T$, and machine $M$.
\end{definition}

\begin{theorem}[Program output encodes instruction bits]
\lean{Complexity.UTab.prog_instr_bits}
\label{Complexity.UTab.prog_instr_bits}
\leanok
\uses{BoolCircuit.BitSrc, Complexity.CounterProg.MOp.exec, Complexity.UTab.ER, Complexity.UTab.ER.f, Complexity.UTab.SSrc.Ok, Complexity.UTab.SSrc.toBitSrc, Complexity.UTab.Tpl, Complexity.UTab.Tpl.kind, Complexity.UTab.Tpl.srcs, Complexity.UTab.ibits, Complexity.UTab.instrOps, Complexity.UTab.instr_bits, Complexity.UTab.tmpl, Complexity.cfgArity, Complexity.cfgInstr, Complexity.cfgProg, Complexity.cfgProg_getElem, Complexity.cfgProg_wf, Complexity.cfgStride, Complexity.cfgWidth, Complexity.length_cfgProg, Complexity.length_tabLayout, Complexity.padSrcs, Complexity.snapWidth, Complexity.tabLayout, Complexity.tabLayout_valid}
\difficulty{4}
Let $M$ be a finite Turing machine over the Boolean alphabet, let $n,T$ be natural
numbers, let $z$ be a Boolean string, let $\theta$ be a template for $M$, and let $e$ be
an eleven-register record. Suppose the template program for $\theta$ equals its
instruction-operation list, register $NC$ in $e$ equals $n$, register $Z$ equals $|z|$,
and register $I$ is less than $(T+1)$ times the configuration stride for input length
$|z|+n$ and time bound $T$. Suppose there is a list $srcs$ of bit sources such that the
configuration instruction at index $e.I$ is the instruction of kind $\theta.kind$ with
its sources padded from $srcs$; the template's sources, translated using the register
values of $e$, equal that padded list; and every template source is valid for those
register values. Then executing the template's operations with the register values of
$e$ produces exactly the bit encoding of that configuration instruction.
\end{theorem}

\begin{theorem}[Rewriting a natural-number equation]
\lean{Complexity.UTab.prev_eq}
\label{Complexity.UTab.prev_eq}
\leanok
\difficulty{3}
Let $S,t,a,q$ be natural numbers with $t\ne 0$. If $q+S=tS+a$, then $q=(t-1)S+a$.
\end{theorem}

\begin{theorem}[Cell Template Bit Output]
\lean{Complexity.UTab.cell_bits}
\label{Complexity.UTab.cell_bits}
\leanok
\uses{Complexity.CounterProg.MOp.exec, Complexity.UTab.ER, Complexity.UTab.ER.f, Complexity.UTab.ER.f_i, Complexity.UTab.ER.f_lb, Complexity.UTab.ER.f_q, Complexity.UTab.Tpl.kind, Complexity.UTab.Tpl.srcs, Complexity.UTab.cellS_toBitSrc, Complexity.UTab.ibits, Complexity.UTab.ok_cellS, Complexity.UTab.ok_padS, Complexity.UTab.padS_toBitSrc, Complexity.UTab.prev_eq, Complexity.UTab.prog_instr_bits, Complexity.UTab.rQ, Complexity.UTab.tmpl, Complexity.cellIdx, Complexity.cellIdx_lt_succ, Complexity.cellSrcs, Complexity.cfgInstr, Complexity.cfgPos_cellIdx, Complexity.cfgStride, Complexity.length_tabLayout, Complexity.padSrcs, Complexity.tabLayout}
\difficulty{5}
Let $M$ be a finite Turing machine over the binary alphabet, and choose natural numbers
$n,|z|,T,t,r$ and a work-tape index $\tau<k$. Suppose $r<2T+1$ and $t\le T$. Let $e$ be
a register state whose input-length and fixed-string-length registers are $n$ and $|z|$,
respectively; whose instruction register is the layout index of cell $(t,\tau,r)$; whose
preceding-cell register, when $t\ne0$, satisfies $q+\mathrm{stride}=i$; and whose
first-instruction register is $t\,\mathrm{stride}$. For a phase $ph$ in the five-element
phase set, suppose its nonzero flag agrees with $1\le r$, its phase-$2$ flag agrees with
$r=T$, and its not-phase-$4$ flag agrees with $r+1<2T+1$. Then executing the cell
template’s operations on $e$ produces exactly the bit string assigned to instruction $i$
by the table construction.
\end{theorem}

\begin{theorem}[Input template matches tableau instruction bits]
\lean{Complexity.UTab.inp_bits}
\label{Complexity.UTab.inp_bits}
\leanok
\uses{Complexity.CounterProg.MOp.exec, Complexity.UTab.ER, Complexity.UTab.ER.f, Complexity.UTab.ER.f_i, Complexity.UTab.ER.f_lb, Complexity.UTab.ER.f_q, Complexity.UTab.IPh, Complexity.UTab.IPh.lay, Complexity.UTab.SSrc.Ok, Complexity.UTab.SSrc.toBitSrc, Complexity.UTab.Tpl.kind, Complexity.UTab.Tpl.srcs, Complexity.UTab.ibits, Complexity.UTab.inpPhS, Complexity.UTab.inpS, Complexity.UTab.inpS_toBitSrc, Complexity.UTab.layS, Complexity.UTab.ok_inpS, Complexity.UTab.ok_padS, Complexity.UTab.padS_toBitSrc, Complexity.UTab.prev_eq, Complexity.UTab.prog_instr_bits, Complexity.UTab.rQ, Complexity.UTab.tmpl, Complexity.cfgInstr, Complexity.cfgPos_inpIdx, Complexity.cfgStride, Complexity.inpIdx, Complexity.inpIdx_lt_succ, Complexity.inpSrcs, Complexity.length_tabLayout, Complexity.padSrcs, Complexity.snapWidth, Complexity.tabLayout}
\difficulty{5}
Let $M$ be a finite-state Turing machine with $k$ work tapes, let $n,T,t,p\in\mathbb N$,
let $z$ be a finite Boolean string, and let $\theta$ be one of the seven input phases.
Suppose $p<|z|+n+2$, $t\le T$, and let $e$ be a register record whose input-count and
hard-wired-input-length registers equal $n$ and $|z|$, whose current-instruction
register is the tableau position for input cell $(t,p)$, and whose step-start register
is $t$ times the configuration stride. If $t\ne0$, assume its preceding-cell register
plus one stride equals its current-instruction register. Assume also that the phase’s
source list agrees with the prescribed input-cell sources, that its layout source
selects the corresponding hard-wired input bit when $1\le p\le |z|+n$ and is false
otherwise, and that this layout source is valid for the register record. Then executing
the template for the input cell of phase $\theta$ produces exactly the bit string of the
corresponding tableau instruction.
\end{theorem}

\begin{theorem}[Snapshot template bit output]
\lean{Complexity.UTab.snap_bits}
\label{Complexity.UTab.snap_bits}
\leanok
\uses{Complexity.CounterProg.MOp.exec, Complexity.UTab.ER, Complexity.UTab.ER.f, Complexity.UTab.ER.f_i, Complexity.UTab.ER.f_lb, Complexity.UTab.Tpl.kind, Complexity.UTab.Tpl.srcs, Complexity.UTab.ibits, Complexity.UTab.ok_padS, Complexity.UTab.ok_snapS, Complexity.UTab.padS_toBitSrc, Complexity.UTab.prog_instr_bits, Complexity.UTab.snapS_toBitSrc, Complexity.UTab.tmpl, Complexity.cfgInstr, Complexity.cfgPos_snapIdx, Complexity.cfgStride, Complexity.length_tabLayout, Complexity.padSrcs, Complexity.snapIdx, Complexity.snapIdx_lt_succ, Complexity.snapSrcs, Complexity.tabLayout}
\difficulty{4}
Let $M$ be a finite Turing machine with Boolean symbols, let $n,T,t$ be natural numbers,
let $z$ be a finite list of bits, and let $e$ be an eleven-register record. Assume $t\le
T$, and that the input-count, hard-wired-length, time-budget, instruction-position, and
step-start registers of $e$ equal, respectively, $n$, the length of $z$, $T$, the index
of the snapshot position at time $t$ for input length $|z|+n$ and budget $T$, and $t$
times the configuration stride for input length $|z|+n$ and budget $T$. The concatenated
output of executing the snapshot template for time $t$ on the register values of $e$
equals the instruction-bit list specified by $M,n,z,T$ at the instruction position
recorded in $e$.
\end{theorem}
